% Clase 4 — Planeación en Ambientes No Deterministas
% Lectura: AIMA Secciones 11.5–11.7
\section{Planeación en Ambientes No Deterministas}

\subsection{Ambientes No Deterministas y Condicionales}

Cuando las acciones tienen efectos inciertos, se necesitan \textbf{planes condicionales} que incluyan ramificaciones según el resultado de las acciones.

\begin{itemize}
    \item \textbf{Plan de búsqueda AND-OR}: los nodos AND representan contingencias del entorno; los nodos OR representan elecciones del agente.
    \item Un plan condicional es una solución si cubre todas las ramas.
\end{itemize}

\subsection{Planeación Parcialmente Observable}

Con observabilidad parcial el agente mantiene un \textbf{estado de creencia} (belief state) y ejecuta acciones de recopilación de información.

\subsection{Planeación Continua y de Ejecución}

\begin{itemize}
    \item \textbf{Monitoreo de ejecución}: detectar discrepancias entre el plan y la realidad.
    \item \textbf{Replaneación}: adaptar el plan ante fallos o nueva información.
    \item \textbf{Agentes reactivos-deliberativos}: combinan planeación a largo plazo con respuesta reactiva.
\end{itemize}

% Agrega aquí ejercicios de planeación condicional
